% !TEX root = ../notes_dgxspark.tex
% Chapter 2: Settle In -- folders, environment, detached runs
%%%%%%%%%%%%%%%%%%%%%%%%%%%%%%%%%%%%%%%%%%%%%%%%%%%%%%%%%%%%%%%%%%%%%%
%
% Sandbox: everything a student does lives under their home directory.
% Sandbag: their runs are niced, time-limited, and detached.
% The venv + cu130 wheel route is the one NVIDIA's own LLaMA Factory
% playbook uses; the NGC container route is in its Unsloth playbook.

\index{settle in}
\chapter{Settle In}
\printchapterversion{02_settlein}

\section{Your Home Is Your Sandbox}

\index{home directory}\newthought{Everything you do lives under your home
directory.} Nothing goes into \texttt{/opt}, \texttt{/usr/local}, or
another user's folder, and nothing you do in this chapter needs
\texttt{sudo}.
\begin{verbatim}
mkdir -p ~/projects ~/data ~/runs
\end{verbatim}
\marginnote{\newthought{Code, data, and results apart.} Code in
\texttt{projects}, datasets in \texttt{data}, logs and checkpoints in
\texttt{runs}. Deleting a run then never touches your code.}

\index{disk}\newthought{The disk is shared too.} Check what you hold, and
delete what you no longer need. Model downloads pile up in
\texttt{\textasciitilde/.cache}.
\begin{verbatim}
du -sh ~ ~/.cache/*
\end{verbatim}

\section{One Environment per Project}

\index{venv}\index{PyTorch}\newthought{A virtual environment keeps your
packages yours.} Create one inside the project and install PyTorch for
CUDA 13 into it:
\begin{verbatim}
cd ~/projects/myproject
python3 -m venv .venv
source .venv/bin/activate
pip install torch --index-url \
  https://download.pytorch.org/whl/cu130
\end{verbatim}
\marginnote{\newthought{Never \texttt{sudo pip install}.} It writes into
the system Python that every other user runs on.}
\marginnote{\index{ARM}\newthought{The Spark is an ARM machine.} A package
without an \texttt{aarch64} wheel builds from source or fails. Check
before you plan around it.}

\section{Start, Detach, Leave}

\index{nohup}\index{nice}\newthought{Run jobs detached, at low priority,
with a time limit.} The job keeps running after you log out, yields the
CPU to interactive users, and stops itself after eight hours:
\begin{verbatim}
nohup nice -n 10 timeout 8h \
  python train.py > ~/runs/train.log 2>&1 &
\end{verbatim}
\marginnote{\newthought{Check on it later} with
\texttt{tail -f \textasciitilde/runs/train.log}. Stop it with
\texttt{kill} and the PID from \texttt{ps -u \$USER}.}
